
\section{Experiments}\label{sec:exp}

\textbf{Datasets.} We evaluate our method on the Voxceleb2~\cite{voxceleb2} and LRW~\cite{lrw} datasets, which contain in-the-wild videos of talking human faces and are  commonly used for lip-sync research. 

\noindent{\textbf{Voxceleb2}}~\cite{voxceleb2} - consists of over 1M face-cropped Youtube videos coming from 6000+ identities. This dataset consists of high variation in lighting, image quality, pose, and motion blur. The average video length is ~8 seconds.

\noindent{\textbf{LRW}}~\cite{lrw} - is a lip-reading dataset that contains 1000 videos each of 500 different words for a length of 1 second coming from BBC news. It has less variation compared to Voxceleb2 and focuses on clean front-facing videos. 
Like previous works, our model is trained only on the Voxceleb2 train split while we test on both datasets. We don't use the whole dataset for training but rather only use the first utterance of every video, which totals 145K videos.  


\textbf{Implementation Details.} We preprocess the videos to have a framerate of 25 fps and an audio sample rate of 16kHz. For all our models the video resolution is 224$\times$224 out of which we crop the face and resize it to 128$\times$128. This is then masked in the lower half using gaussian noise and fed to our model which only morphs the lower half of this image according to the audio input. Then we resize it back to the original crop size and place it back on the video. For audio inputs, we first sample the audio at 16kHz and then create mel-spectrograms with window-size 800 and hop-size 200. These audio frames turn out to have size 16$\times$80. We build our code on top of the guided-diffusion repository~\cite{guided-diffusion}. %
We train our model on 8 NVIDIA RTXA6000 GPUs which takes around 4 days. Our model is trained using $T=1000$ diffusion steps, but for faster inference, we use only 25 steps of DDIM~\cite{ddim} sampling which takes 4.67 seconds on an average for all the frames of one VoxCeleb2~\cite{voxceleb2} video (avg. 8 seconds at 25 fps) on 8 NVIDIA RTXA6000 GPUs.


\begin{table}[!t]
\setlength{\cmidrulewidth}{0.01em}
\renewcommand{\tabcolsep}{3pt}
\renewcommand{\arraystretch}{1.1}
\caption{Ablation over our losses (Reconstruction)}\label{tab:ablation_recon}
\centering
\resizebox{\linewidth}{!}{
\begin{tabular}{@{}lcccccc@{}}
\toprule
Losses & FID ↓ & SSIM ↑& PSNR ↑ & LMD ↓ & Sync$_\text{c}$ ↑  \\
 \midrule
Reconstruction          & 8.589          & 0.523          & 18.234         & 3.472          & 0.633         \\
+ SyncNet     & 8.998          & 0.526 & \textbf{18.57} & 3.123          & 6.336         \\
+ Perceptual & \textbf{7.751} & 0.526          & 18.548         & 3.121          & 6.53          \\
+ Seq. GAN           & 8.213 & \textbf{0.527} & 18.52 & \textbf{3.101} & \textbf{7.89} \\
\bottomrule
\end{tabular}
}
    \vspace{-0.2in}
\end{table}


\begin{table*}[!t]
\centering
\small
\setlength{\cmidrulewidth}{0.01em}
\renewcommand{\tabcolsep}{4pt}
\renewcommand{\arraystretch}{1.1}
\caption{Quantitative comparison with baselines on Voxceleb2~\cite{voxceleb2} and LRW~\cite{lrw} on the task of reconstruction and Cross generation.}
\label{tab:main_table}
\begin{tabular}{@{}llcccccccccc@{}}
\toprule
\multirow{2}{*}{Dataset} &
\multirow{2}{*}{Method} & \multicolumn{6}{c}{Reconstruction} & \multicolumn{4}{c}{Cross} \\
\cmidrule[\cmidrulewidth](l){3-8} \cmidrule[\cmidrulewidth](l){9-12}
& & FID ↓ & SSIM ↑& PSNR ↑ & LMD ↓ & Sync$_\text{c}$ ↑ & Sync$_\text{d}$ ↓ & FID ↓ & LMD ↓ & Sync$_\text{c}$ ↑ & Sync$_\text{d}$ ↓ \\
 \midrule

\multirow{3}{*}{VoxCeleb2} & Wav2Lip~\cite{wav2lip} & 3.26 & 0.53   & 18.18                        &  3.16 & \textbf{9.08} & 5.93               & 5.11               & 4.84               & \textbf{8.12} & 6.74              \\ 
 & PC-AVS~\cite{pcavs} & 4.25                        & 0.53               & \textbf{18.26}               & 3.16                        &  6.71  & 7.80 &  10.62 &  5.00 &  6.96 & 7.53\\ 
& \ours~(Ours) & \textbf{2.46}               & 0.53 &  18.09 & \textbf{3.04}               & 8.78           & \textbf{5.93}             & \textbf{4.53}                & \textbf{4.82}               &  7.62  & \textbf{6.73}  \\ 

\midrule
\multirow{3}{*}{LRW} & 
Wav2Lip~\cite{wav2lip}  &  4.23 & \textbf{0.68} & \textbf{20.76} & \textbf{2.15} & \textbf{8.13}  & \textbf{6.09 }         & 5.19 & \textbf{3.88} & \textbf{7.52} & \textbf{6.56}\\ 
 & PC-AVS~\cite{pcavs}     & 6.80 & 0.61 & 20.10 & 2.29 & 6.68 & 7.29
 & 8.48 & 4.09 & 6.66 & 7.27 \\ 
 & \ours~(Ours)            & \textbf{2.62} & 0.67 & 20.62 & 2.17 & 7.41 & 6.21
 & \textbf{2.54} & 3.93 & 6.44 & 6.97\\ 
\bottomrule
\end{tabular}

    \vspace{-0.1in}
\end{table*}






\textbf{Comparison Methods.} We compare our method against the most popular methods for lip-sync. Our choice of models is based also on models/codebases which are publicly available. Wav2Lip~\cite{wav2lip} is an inpainting style method that uses SyncNet expert loss to get good lip-sync. PC-AVS~\cite{pcavs} is a head reconstruction method that focuses on controlling pose apart from identity and lip shape. For both these methods we use their publicly available pre-trained models for the evaluation of all the datasets.
















\begin{table}[!t]
\begin{minipage}[t]{\linewidth}
\setlength{\cmidrulewidth}{0.01em}
\renewcommand{\tabcolsep}{9pt}
\renewcommand{\arraystretch}{1.1}
\caption{Ablation over our losses (Cross generation)}\label{tab:ablation}
\centering
\resizebox{0.8\linewidth}{!}{
\begin{tabular}{@{}lcccc@{}}
\toprule
Losses & FID ↓ & LMD ↓ & Sync$_\text{c}$ ↑ \\
 \midrule

Reconstruction          & 6.694          & 5.313          & 0.992         \\
+ SyncNet     & 8.784          & \textbf{4.816} & 5.946         \\
+ Perceptual & 5.016          & 5.009          & 5.955          \\
+ Seq. GAN           & \textbf{4.592} & 4.985          & \textbf{6.83} \\
\bottomrule
\end{tabular}
}
\end{minipage}

    \vspace{-0.1in}
\end{table}



\begin{figure*}[!t]
    \centering
    \label{fig:qual_recon}
    \includegraphics[width=0.9\linewidth]{figures/qualitative-reconv2.pdf}
    \caption{\textbf{Qualitative results of Reconstruction on VoxCeleb2~\cite{voxceleb2}}. Here we provide only the first frame as the input source (for pose) as well as the reference frame (for identity), and this frame is driven using the audio (second row) coming from the same video (top row). Wav2Lip~\cite{wav2lip} blurs the lip region in both cases to achieve the correct lip shape while PC-AVS has identity loss (see right) and border discontinuity. Our generations look highly realistic and have the lip shapes as in the audio source. (Please zoom in for better visibility.)}
    \label{fig:qualitative-recon}
    \vspace{-0.1in}
\end{figure*}

\begin{figure*}[!h]
    \centering
    \label{fig:qual_cross}
    \includegraphics[width=0.9\linewidth]{figures/qualitative-crossv2.pdf}
    \caption{\textbf{Qualitative results of Cross generation on VoxCeleb2~\cite{voxceleb2}}. Here we provide a video source (first row) and drive that identity-pose combination using audio coming from a different video (second and third rows). Wav2Lip~\cite{wav2lip} blurs its generations, for example, beard region details are missing on the right. PC-AVS's~\cite{pcavs} generations have flaws like identity loss, in both cases. They introduce artifacts near the eyes on the left while there is identity loss on the right. Our method generates realistic mouths with expressive lips while being in sync with the audio source. In the bottom 2 rows, we can see that the lip region of our generations match those of the audio source. (Please zoom in for better visibility.)}
    \label{fig:qualitative-cross}
    \vspace{-0.1in}
\end{figure*}

\subsection{Quantitative Evaluation}\label{sec:quant-eval}
 For quantitative evaluation, we evaluate our model in terms of both the visual quality as well as audio synchronization. For visual quality, we use FID~\cite{fid}, SSIM~\cite{ssim}, and PSNR, which are popular metrics used in papers like Wav2Lip~\cite{wav2lip}, PC-AVS~\cite{pcavs}, AV-CAT~\cite{avcat}. FID is a popular metric used for comparing the ``realness'' of generated images by comparing against the real image distribution. SSIM and PSNR are pixel-wise image similarity metrics that compare a pair of images and are not suited for capturing variability in video generation~\cite{shrivastava2021diverse} but are included in this work for completeness. While to measure synchronization we use LMD~\cite{chen2018lip}, $\text{Sync}_\text{c}$, and $\text{Sync}_\text{d}$~\cite{syncnet}. LMD measures the distance between mouth landmarks among frames. $\text{Sync}_{c}$ is the confidence score of SyncNet while $\text{Sync}_\text{d}$ is the average distance between SyncNet video and audio representations, which tell the synchronization quality. Note that for evaluation, we use the pre-trained SyncNet from the SyncNet's~\cite{syncnet} repository but for training, we train our own SyncNet, similar to Wav2Lip\cite{wav2lip} and AV-CAT\cite{avcat}.
We use pre-trained models of Wav2Lip and PC-AVS methods to conduct our evaluations. Wav2Lip has provided its code for calculating FID and $\text{Sync}_\text{c}$ and we use the same for these metrics. We use face-alignment ~\cite{face-alignment} for landmark detection and the LMD metric proposed in~\cite{chen2018lip}. For SSIM and PSNR we use the same inputs as used for FID, consequently, our values are a bit different compared to~\cite{pcavs}. This is  possibly because they evaluate these metrics at different scales and lack these evaluation details. On the other hand, we tend to get better LMD and $\text{Sync}_\text{c}$ values for PC-AVS than noted in their paper. Some papers like~\cite{gcavt} show PC-AVS's metrics only for reference as its generations occasionally fail in landmark detection. For uniformity of scale, we uncropped PC-AVS's generations and paste them back on the background before evaluation.  




\begin{figure*}[!t]
    \centering
    \includegraphics[width=0.95\linewidth]{figures/qualitative-zoomed-y.pdf}
    \caption{\textbf{Qualitative Visual-quality Comparison}. We zoom in on the mouth region of two examples and compare them against the video source. Wav2Lip~\cite{wav2lip} blurs the lip region while PC-AVS~\cite{pcavs} tends to change the identity. \ours~ preserves identity and generates high-fidelity lips.}
    \label{fig:qualitative-zoomed}
    \vspace{-0.1in}
\end{figure*}


\subsubsection{Reconstruction}\label{sec:recon}
Similar to the setting mentioned in~\cite{pcavs} and later used in~\cite{gcavt} and~\cite{avcat}, we evaluate the methods on the task of reconstruction of the video given only the first frame and the audio corresponding to the same video. The audio is used as the driver for this reconstruction. Note that PCAVS requires an additional pose input apart from the identity input. We feed the first frame to it for both inputs, similar to the ``fixed pose'' setting in their paper. For Wav2lip and ours, the first frame acts as both the input frame and the reference frame. For Voxceleb2, we show the results on 4911 test utterances instead of 35K videos due to resource constraints. These 4911 utterances are the first utterance of each video in the test set, and hence cover all the videos in the test set. For LRW, our results are noted on all 25K videos in the test set.

The results are noted in Table~\ref{tab:main_table}. We observe that \ours~outperforms both the methods with respect to FID metric on both datasets which points towards better generation quality. The SSIM, PSNR and LMD values of our method are comparable with the other methods. We see that Wav2Lip's  $\text{Sync}_\text{c}$ is better than both ours and PC-AVS's. This is possibly because our SyncNet expert may be weaker in performance compared to the one used in Wav2Lip. 

\vspace{-0.1in}
\subsubsection{Cross generation}\label{sec:cross}
\vspace{-0.1in}
We also evaluate the method on the task of lip-sync when the identity and the pose are controlled using a video while the lip-sync is driven using input audio corresponding to a different video. This was introduced in~\cite{wav2lip} and is a more realistic setting as here the generations are closer to lip-sync in-the-wild. We use the input frame as the reference frame in this setting similar to Wav2Lip, as that provides the best texture information of the frame. For PC-AVS, the input frames are fed as both pose and identity sources. Note that we cannot evaluate SSIM and PSNR for this setting because there are no ground truth frames available. So, we provide the rest of the metrics in Table~\ref{tab:main_table}. For Voxceleb2, we select 4970 pairs of audio-video combinations where the two sources are different. We sample these using all the pair combinations of the first utterance of the first video coming from 71 randomly chosen test identities. For LRW, we use the 28K audio-video pair provided in Wav2Lip's~\cite{wav2lip} evaluation.  
 The results are noted in Table~\ref{tab:main_table}. Similar to reconstruction evaluation, we here as well observe that our method excels in image quality while being comparable in other metrics except for $\text{Sync}_\text{c}$.

\subsection{Qualitative Evaluation}
For qualitative evaluation, we show visually compare against on both reconstructions (in Fig.~\ref{fig:qualitative-recon}) as well as cross generation (in Fig.~\ref{fig:qualitative-cross}). These settings are the same as introduced in Section~\ref{sec:recon} and \ref{sec:cross}. It can be observed in these qualitative results that PC-AVS tends to lose the identity of the source video and also suffers from boundary discontinuity problems which make it unsuitable for in-the-wild generation. On the other hand, Wav2Lip tends to generate blurred-out mouth regions so as to achieve good lip sync. \ours~ does not suffer from these issues and is able to generate high-quality mouth region while having expressive lip shapes which correctly correspond to the ground truth (audio sources' mouth shape) as seen in Fig.~\ref{fig:qualitative-zoomed}.




\begin{table}[!t]
\small
\setlength{\cmidrulewidth}{0.01em}
\renewcommand{\tabcolsep}{3pt}
\renewcommand{\arraystretch}{1.1}
\caption{User Study measured by Mean Opinion Scores~(MOS) (max. 5) and  Preference in percentage.}\label{tab:user_study}
\centering
\begin{tabular}{@{}lccc@{}}
\toprule
Measure & Wav2Lip & PC-AVS  & \ours \\
 \midrule

MOS (Visual quality) ↑   & 3.75 & 2.71 & \textbf{4.16} \\
MOS (Lip-sync quality) ↑ & 3.84 & 3.34 & \textbf{3.86} \\
MOS (Overall quality) ↑  & 3.70 & 2.91 & \textbf{3.91} \\
Preference ↑         & 37\% & 8.33\% & \textbf{54.67}\% \\
\bottomrule
\end{tabular}
    \vspace{-0.1in}
\end{table}

\textbf{User Study.} We conduct a user study where we ask 15 participants to judge lip-sync videos generated in cross generation setting by \ours~ and two other methods. 20 videos were sampled from the VoxCeleb2's test set and are driven by randomly selected driving audios. The participants rated the videos 1-5 (where higher is better) in the aspects of (1) \textbf{Visual quality} (2) \textbf{Lip-sync quality}, and (3) \textbf{Overall quality}. We used the Mean Opinion Score (MOS) measure to aggregate these ratings. Further, we record the percentage of times users preferred a method. We present the results in Table~\ref{tab:user_study}, where we see that our method surpasses others in all the categories. In terms of Lip-sync quality, this is opposed to our quantitative results, especially $\text{Sync}_\text{c}$. We speculate that $\text{Sync}_\text{c}$, might favor blurry generations with high temporal consistency while humans prefer high fidelity over slight temporal inconsistency. 



\subsection{Ablations}
We conduct an ablation study to showcase the contribution of various choices of losses that were used during training. Specifically, we train our model in three different settings in which we introduce an additional loss in each setting. First, in the setting, we train our model using only $\mathcal{L}_\text{simple}$ (Reconstruction). Second, we train another version of the model 
 using $\mathcal{L}_\text{simple}+\mathcal{L}_{2}+\mathcal{L}_\text{sync}$ (+ SyncNet). Here intuitively the $\mathcal{L}_\text{sync}$ should introduce better synchronization. Third, we further add a perceptual loss $\mathcal{L}_\text{sync}$ (+ Pecept). We add this loss because adding the SyncNet loss led to worse image quality. Finally, we add the sequential adversarial loss $\mathcal{L}_\text{GAN}$ to achieve even temporal consistency(+ Seq. GAN). We test these on a smaller subset of 500 VoxCeleb2 test audio-video pairs in the cross generation setting as well as the reconstruction setting. 
 
It can be seen in Table~\ref{tab:ablation} that moving from ``Reconstruction" to ``+ SyncNet" gives a sudden improvement in the $\text{Sync}_\text{c}$ metric. This supports our intuition that only reconstruction-based losses are not enough. We also see that this transition deteriorates the image quality. This gets solved as we move to the ``+ Perceptual" setting. Finally, the addition of sequential adversarial loss not just further improves the image quality but also improves the $\text{Sync}_c$, clearly showing the advantage of this loss.  In the reconstruction setting in Table~\ref{tab:ablation_recon}, most of these observations still hold except FID being lower for ``+ Perceptual" than ``+ Seq. GAN". This could be attributed to the static nature of the input source in this setting while the ground truth is moving.  


